\documentclass[11pt]{article}
\usepackage[T1]{fontenc}
\usepackage{lmodern}
\usepackage{graphicx}
\usepackage{caption}
\usepackage{subcaption}
\usepackage{booktabs}
\usepackage{multirow}
\usepackage[table]{xcolor}
\usepackage{natbib}
\usepackage{hyperref}
\begin{document}
\title{A Report on Memory Hierarchies}
\author{Henry Yang}
\date{\today}
\maketitle
\tableofcontents

\section{Introduction}
Memory hierarchies trade capacity for speed \citep{hennessy2017}. \citet{smith1982} surveyed cache designs; later work \cite{hennessy2017,smith1982} refined them. Figure~\ref{fig:both} shows two placeholder images.

\begin{figure}[h]
  \centering
  \begin{subfigure}[b]{0.45\textwidth}
    \includegraphics[width=\linewidth]{images/figure.png}
    \caption{First.}\label{fig:a}
  \end{subfigure}
  \hfill
  \begin{subfigure}[b]{0.45\textwidth}
    \includegraphics[width=\linewidth]{images/figure.png}
    \caption{Second.}\label{fig:b}
  \end{subfigure}
  \caption{Two placeholder images.}\label{fig:both}
\end{figure}

\section{Results}
Table~\ref{tab:res} summarises latencies.

\begin{table}[h]
\centering
\rowcolors{2}{gray!15}{white}
\begin{tabular}{llrr}
\toprule
\multirow{2}{*}{Level} & \multirow{2}{*}{Technology} & \multicolumn{2}{c}{Latency} \\
\cmidrule(lr){3-4}
 & & ns & cycles \\
\midrule
L1 & SRAM & 1 & 4 \\
L2 & SRAM & 4 & 12 \\
\rowcolor{yellow!30} DRAM & capacitors & 60 & 200 \\
\bottomrule
\end{tabular}
\caption{Typical latencies.}\label{tab:res}
\end{table}

As \citeauthor{smith1982} noted in \citeyear{smith1982}, associativity matters. The data in Table~\ref{tab:res} are illustrative only.

\section{Conclusion}
Caches work because of \textbf{locality}. See \autoref{fig:a} and page \pageref{tab:res}.

\bibliographystyle{plainnat}
\bibliography{refs}
\end{document}
